\subsubsection{La sala de ejemplo}\label{sec:sala}

Igual que con el casco, cada paso de la reconstrucción se resuelve a continuación, dentro del mismo paso, con una fotografía real y con los números que produce el programa. La
foto es una sala con un sofá, una mesa de centro y un mueble de televisión (figura~\ref{fig:hab-entrada}), tomada de Wikimedia Commons
(autor: Shixart1985, licencia CC BY 2.0, \texttt{commons.wikimedia.org}). Se reduce a $960\times640$ píxeles ($N_p=614\,400$ píxeles) y se procesa con el
programa tal como está (procesamiento bajo); cada etapa se ejecutó por separado
(\texttt{experimentos/habitacion\_paso\_a\_paso.py}). En la figura~\ref{fig:hab-flujo} están los seis pasos y los objetos matemáticos de cada uno.

\begin{figure}[htbp]
\centering
\includegraphics[width=0.95\linewidth]{hab_entrada.jpg}
\caption{Foto de la sala reducida a $960\times640$ píxeles, con los ejes $u$ (columna) y $v$ (fila), el centro $(c_x,c_y)$ y los cuatro píxeles
(piso, pared, techo y sofá) que se siguen en las cuentas.}
\label{fig:hab-entrada}
\end{figure}

\begin{figure}[htbp]
\centering
\includegraphics[width=\linewidth]{hab_flujo.jpg}
\caption{Los seis pasos de la reconstrucción de una habitación y el objeto matemático de cada uno.}
\label{fig:hab-flujo}
\end{figure}


\paragraph{La cámara de la sala.} Con el ángulo de visión $\alpha=65^\circ$ y $W=960$, la distancia focal es $f=(W/2)/\tan(\alpha/2)=753{,}449$ píxeles,
$(c_x,c_y)=(480,320)$ y
\begin{equation}
K=\begin{pmatrix} 753{,}449 & 0 & 480 \\ 0 & 753{,}449 & 320 \\ 0 & 0 & 1 \end{pmatrix},\qquad \det K=f^2=567\,685{,}5,
\end{equation}
\begin{equation}
K^{-1}=\begin{pmatrix} 0{,}001327 & 0 & -0{,}637070 \\ 0 & 0{,}001327 & -0{,}424714 \\ 0 & 0 & 1 \end{pmatrix} .
\end{equation}
Son los mismos valores que en el casco porque la foto tiene el mismo ancho; solo cambia $c_y$ porque la altura es otra.
